\chapter{Сдвинутая кривизна: формульный допуск и градуировочный запрет}
\label{ch:v8-baryon-spacetime-supercurvature-cubic-projection-admission-gate}

Предыдущая глава оставила один допустимый маршрут: кубический член должен
возникнуть вместе с квадратичным и четвертичным членами из единой нормы
кривизны. База дедукции подсказывает квадрат отображения момента, а
суперсвязностная литература --- постоянную нечётную часть конечного
оператора Дирака. Проверим их общий минимальный полином на полном
центрированном кадре.

\section{Единая сдвинутая кривизна}

Пусть
\begin{equation}
 Z=z^a\widehat F_a,\qquad
 M=\frac m2I_{21}.
 \label{eq:v8-baryon-shifted-supercurvature-field}
\end{equation}
Центральность $M$ даёт ковариантную при сопряжении относительную кривизну
\begin{equation}
 \mathcal R_m(Z)=(M+Z)^2-M^2=mZ+Z^2.
 \label{eq:v8-baryon-shifted-supercurvature}
\end{equation}
Алгебраически это точный аналог сдвинутого квадрата отображения момента.
Чтобы считать его конечномерной частью суперсвязности, необходимо отдельно
проверить полную градуировку.

Одна положительная следовая норма раскрывается точно:
\begin{equation}
 \Tr\mathcal R_m(Z)^2
 =m^2\Tr Z^2+2m\Tr Z^3+\Tr Z^4.
 \label{eq:v8-baryon-shifted-supercurvature-action}
\end{equation}
Следовательно, при общей нормировке
\begin{equation}
 \alpha=m^2,\qquad \lambda_3=2m,\qquad \beta=1,
 \qquad \frac{\lambda_3^2}{\alpha\beta}=4.
 \label{eq:v8-baryon-shifted-supercurvature-shape}
\end{equation}
В отличие от предыдущего произвольного полинома, безразмерная форма
коэффициентов теперь фиксирована.

\section{Проекция на связный тензор}

Для $Z=z^a\widehat F_a$ имеем
\begin{equation}
 \Tr Z^3=d_{abc}z^az^bz^c.
 \label{eq:v8-baryon-cubic-trace-polynomial}
\end{equation}
Поэтому кубическая часть единой нормы равна
\begin{equation}
 S^{(3)}(Z)=2m\,d_{abc}z^az^bz^c,
 \qquad
 \delta^3S\big|_0=12m\,d_{abc}.
 \label{eq:v8-baryon-supercurvature-cubic-projection}
\end{equation}
Она имеет уже вычисленную опору $140\,TTG+28\,GGG$ и нулевые каналы
$TTT,TGG$. Значит, единая кривизна воспроизводит не только степень, но и
точную тензорную форму канонического оператора $W_3$.

\section{Точный радиальный контроль}

На прежнем луче $Z=t(\widehat F_0+\widehat F_{40})$ получаем
\begin{equation}
 S_m(t)=38m^2t^2-6mt^3+134t^4.
 \label{eq:v8-baryon-shifted-supercurvature-ray-action}
\end{equation}
Его производная равна
\begin{equation}
 S_m'(t)=2t\left(38m^2-9mt+268t^2\right),
 \qquad
 \Delta=-40655m^2.
 \label{eq:v8-baryon-shifted-supercurvature-ray-stationarity}
\end{equation}
При вещественном $m\ne0$ ненулевой стационарной точки на этом луче нет, а
$S_m''(0)=76m^2>0$. Следовательно, найденная форма задаёт согласованное
взаимодействие, но сама не создаёт барионный вакуум на проверенном срезе.

\begin{theorem}[Формульный допуск кубического коэффициента]
Единая положительная норма сдвинутой кривизны канонически порождает
квадратичный, кубический и четвертичный следовые инварианты и фиксирует
безразмерное отношение $\lambda_3^2/(\alpha\beta)=4$. Её кубическая
вариация пропорциональна ранее построенному связному тензору $d_{abc}$.
\end{theorem}

\section{Граница происхождения}

Все $\widehat F_a$ бесследовы, тогда как
\begin{equation}
 \Tr M=\frac{21}{2}m.
 \label{eq:v8-baryon-central-background-trace}
\end{equation}
Поэтому ненулевой центральный фон не лежит в текущем 42-мерном
центрированном пространстве флуктуаций. Ни его масштаб $m$, ни полная
пространственно-временная степень формы, ни оператор Ходжа пока не выведены
существующим родителем. Литературный механизм доказывает допустимость
формы, но не происхождение конкретного $m$ в данной модели.

Более того, для конечной градуировки
$\Gamma=I_{11}\oplus(-I_{10})$ постоянная часть оператора Дирака должна
быть нечётной. Центральный кандидат имеет точный дефект
\begin{equation}
 \Gamma M+M\Gamma=m\Gamma.
 \label{eq:v8-baryon-central-background-parity-defect}
\end{equation}
Он нечётен только при $m=0$, когда кубический член исчезает. Полный кадр
также смешивает конечные чётные калибровочные направления и нечётные
направления переноса; однородность полной суперсвязности восстанавливается
только с учётом пространственно-временной степени дифференциальных форм.
Следовательно, конечный полином сам по себе ещё не является настоящей
суперкривизной.

\begin{nogocriterion}[Запрет подмены формы числом]
Равенство $\lambda_3=2m$ нельзя объявлять вычислением физического
коэффициента или суперсвязностным происхождением. Центральный ненулевой фон
уже запрещён нечётностью. Следующий гейт должен классифицировать
нецентральные нечётные фоновые операторы и проверить, может ли их кубическая
проекция быть пропорциональна всему $d_{abc}$, а не только отдельному
сектору.
\end{nogocriterion}

\begin{protocol}
\begin{enumerate}
 \item поле вычислений: \textbf{$\mathbb Q(m,t)$};
 \item числа с плавающей точкой: \textbf{отсутствуют};
 \item единая положительная норма: \textbf{получена};
 \item кубическая проекция на $d_{abc}$: \textbf{$2m\,d_{abc}$};
 \item безразмерное отношение формы: \textbf{$4$};
 \item ненулевой радиальный стационарный пункт: \textbf{отсутствует};
 \item центральный фон лежит в кадре: \textbf{нет};
 \item центральный фон нечётен по $\Gamma$: \textbf{нет при $m\ne0$};
 \item масштаб $m$ выведен: \textbf{нет};
 \item полная пространственно-временная норма: \textbf{ещё не построена};
 \item формульная форма связного взаимодействия: \textbf{допущена};
 \item настоящая градуированная суперкривизна: \textbf{не получена};
 \item физический коэффициент: \textbf{не выведен}.
\end{enumerate}
\end{protocol}

Машинный сертификат сохранён в
\path{s2t_v8_baryon_spacetime_supercurvature_cubic_projection_admission_gate_results.json}.

\begin{thebibliography}{9}
\bibitem{v8-baryon-supercurvature-roepstorff}
G.~Roepstorff,
\emph{Superconnections and the Higgs Field},
arXiv:hep-th/9801040.
\bibitem{v8-baryon-supercurvature-lee-neeman}
C.-Y.~Lee, Y.~Ne'eman,
\emph{Massive Vector Gauge Theory and Comparison with Higgs--Connes--Lott Theory},
arXiv:hep-th/9807089.
\end{thebibliography}